\begin{circuitikz}[scale=0.85, every node/.style={transform shape}, font=\sffamily\small]
% ---------------- POW-BB rail splitter
\draw[rounded corners=2pt, fill=hwcol] (0,0.6) rectangle (3.2,7.0);
\node[font=\sffamily\bfseries\small] at (1.6,6.5) {SBC-POW-BB};
\node[font=\sffamily\scriptsize] at (1.6,6.1) {labelled rail splitter};
\node[font=\sffamily\scriptsize, align=center] at (1.6,3.0)
  {3V3 comes from\\the board regulator,\\not a guess};
\draw (1.0,7.0) -- (1.0,7.6);
\node[anchor=south, font=\sffamily\scriptsize] at (1.0,7.65) {PSU 5 V in};
% ---------------- rail stubs
\draw (3.2,5.6) -- (9.0,5.6);
\node[anchor=east, font=\sffamily\scriptsize] at (3.05,5.6) {5V};
\node[anchor=south, font=\sffamily\scriptsize] at (5.6,5.75) {5 V rail, labelled: measure 5 V here};
\draw[red!70!black, dashed] (9.0,5.6) -- (10.6,5.6);
\draw[red!70!black, line width=1pt] (9.55,5.35) -- (10.05,5.85);
\draw[red!70!black, line width=1pt] (9.55,5.85) -- (10.05,5.35);
\draw[red!70!black, line width=1pt] (10.6,5.35) -- (10.6,5.85);
\node[anchor=south, font=\sffamily\scriptsize, text=red!70!black, align=center] at (9.8,5.95)
  {no 5 V onto a header pin};
\draw (3.2,4.4) -- (4.6,4.4);
\node[anchor=east, font=\sffamily\scriptsize] at (3.05,4.4) {3V3};
\draw (4.6,4.4) -- (4.6,2.0);
\node[anchor=east, font=\sffamily\scriptsize] at (3.05,1.0) {GND};
% ---------------- three LED branches: anode on 3V3, cathode into a GPIO
\draw (4.6,4.4) to[R, l=series R, -] (6.8,4.4) to[led, -] (8.6,4.4) -- (11.6,4.4);
\draw (4.6,3.2) to[R, l=series R, -] (6.8,3.2) to[led, -] (8.6,3.2) -- (11.6,3.2);
\draw (4.6,2.0) to[R, l=series R, -] (6.8,2.0) to[led, -] (8.6,2.0) -- (11.6,2.0);
\node[anchor=south, font=\sffamily\scriptsize] at (10.0,4.48) {green};
\node[anchor=south, font=\sffamily\scriptsize] at (10.0,3.28) {yellow};
\node[anchor=south, font=\sffamily\scriptsize] at (10.0,2.08) {red};
% ---------------- host
\draw[rounded corners=2pt, fill=hwcol] (11.6,1.3) rectangle (15.4,6.6);
\node[font=\sffamily\bfseries\small] at (13.5,6.2) {Raspberry Pi or NanoPi};
\node[font=\sffamily\scriptsize] at (13.5,5.8) {every I/O pin is 3.3 V};
\node[anchor=west, font=\sffamily\scriptsize] at (11.75,4.4) {pin 13\ \ GPIO27};
\node[anchor=west, font=\sffamily\scriptsize] at (11.75,3.2) {pin 15\ \ GPIO22};
\node[anchor=west, font=\sffamily\scriptsize] at (11.75,2.0) {pin 16\ \ GPIO23};
\node[anchor=west, font=\sffamily\scriptsize] at (11.75,1.58) {a low output sinks the LED};
% ---------------- common ground
\draw (3.2,1.0) -- (14.8,1.0);
\draw (14.8,1.0) -- (14.8,1.3);
\node[anchor=east, font=\sffamily\scriptsize] at (14.7,0.75) {pin 6 GND};
\node[anchor=north, font=\sffamily\scriptsize] at (7.0,0.9) {common ground: POW-BB, breadboard and host};
\draw (10.6,1.0) -- (10.6,0.6);
\node[ground] at (10.6,0.6) {};
% ---------------- negative test, a separate fragment
\draw[dashed, rounded corners=3pt, draw=red!60!black, fill=red!3] (0,-5.0) rectangle (15.4,-0.4);
\node[anchor=north west, font=\sffamily\bfseries\small, text=red!60!black] at (0.2,-0.6)
  {Negative test: on the breadboard only, never on a GPIO};
\node[circle, draw=linecol, fill=white, minimum size=2.2mm, inner sep=0pt] at (1.8,-2.4) {};
\node[anchor=east, font=\sffamily\scriptsize] at (1.6,-2.4) {POW-BB 5 V};
\draw (1.8,-2.4) to[R, l=pull-up R, -] (4.6,-2.4) -- (6.2,-2.4);
\fill (6.2,-2.4) circle (1.3pt);
\node[anchor=south, font=\sffamily\scriptsize] at (6.2,-2.25) {test node};
\draw (6.2,-2.4) -- (6.2,-3.1);
\node[circle, draw=linecol, fill=white, minimum size=7mm, font=\sffamily\small] (vm) at (6.2,-3.45) {V};
\draw (vm.south) -- (6.2,-4.2);
\node[ground] at (6.2,-4.2) {};
\node[anchor=west, font=\sffamily\scriptsize, align=left] at (7.0,-3.6)
  {the meter reads about 5 V,\\above the 3.3 V input limit};
\draw[red!70!black, dashed] (6.2,-2.4) -- (10.8,-2.4);
\draw[red!70!black, line width=1pt] (8.25,-2.65) -- (8.75,-2.15);
\draw[red!70!black, line width=1pt] (8.25,-2.15) -- (8.75,-2.65);
\node[anchor=south, font=\sffamily\scriptsize, text=red!70!black] at (8.5,-2.2) {never};
\node[draw=black!60, fill=black!8, minimum width=24mm, minimum height=9mm, align=center,
      font=\sffamily\scriptsize] at (12.2,-2.4) {a Pi or NanoPi GPIO\\3.3 V maximum};
\node[anchor=west, font=\sffamily\scriptsize, text=red!60!black] at (7.0,-4.4)
  {take the pull-up apart when it has been read};
\end{circuitikz}
